\section{Introduction}
\label{sec:intro}


Memory has become a core component of modern Large Language Model (LLM) agents, enabling them to retain and reuse information beyond a single context window for long-horizon interaction, personalization, and knowledge-intensive reasoning~\citep{hu2025memory}. 
Nevertheless, most prior work centers on offline, query-agnostic memory construction, where past context is preprocessed, compressed, or indexed in a fixed manner without conditioning on the downstream query~\citep{kang2025memory, fang2025lightmem, zhong2024memorybank, xu2025mem, packer2023memgpt, chhikara2025mem0}.
This ``\textit{build once, use always}'' paradigm can be wasteful and brittle: it spends computation regardless of what a particular query needs, while potentially omitting crucial information for specific queries.




Instead, an intuitive alternative is on-demand memory extraction, where computation is triggered at runtime based on the current query.
This flexibility comes at a cost: it pushes memory processing to runtime, making cost and latency first-class concerns.
In practice, industrial LLM systems increasingly provide explicit, often \textit{tiered} compute controls (e.g., ``thinking'' modes~\citep{singh2025openai}, reasoning levels~\citep{openai-codex}, or heavier-model options~\citep{claude-code}), reflecting the need to balance quality against runtime cost.
This motivates a key question for agent memory: \textit{how can we enable explicit and controllable performance–cost trade-offs for runtime memory extraction?}



Unfortunately, enabling performance-cost trade-offs for runtime agent memory is fundamentally challenging. Most existing trade-off mechanisms operate offline, while on-demand memory pushes these decisions to runtime, where each query raises a quality-cost choice. This exposes two core questions.
First, \textit{where should budgets be applied?} Existing systems often treat memory as a monolithic pipeline with a fixed compute setting, making trade-offs coarse and difficult to control~\citep{kang2025memory, zhong2024memorybank}. 
This motivates defining an appropriate budgeting unit for runtime settings, i.e., which modular part(s) of the memory extraction process should be assigned budgets, so that computation can be controlled in a targeted and effective way.
Second, \textit{how should budgets be realized?} Prior work provides little systematic guidance on trade-offs for runtime memory, often resorting to ad hoc cost reduction or simply increasing compute~\citep{yan2025general}. As a result, even after a budgeting scheme is specified, it remains unclear how to operationalize budget control, which design axes best capture meaningful trade-offs, and how these choices behave across different budget regimes.
Addressing these questions requires moving beyond one-off heuristics or compute-heavy escalation toward a more systematic view of performance–cost control for runtime agent memory.


To address the above challenges, we propose \textbf{BudgetMem}, a runtime agent memory framework that enables explicit, controllable performance–cost trade-offs for on-demand memory extraction. 
BudgetMem views runtime memory extraction as a multi-stage modular pipeline and makes computation controllable at the module level. 
Specifically, BudgetMem standardizes how each module is invoked by exposing a common budget-tier interface, so that a learned router can select among budget tiers within modules while keeping the overall extraction structure fixed.


Building on this modular backbone, BudgetMem provides three budget tiers (i.e., \textsc{Low}/\textsc{Mid}/\textsc{High}) for each module, offering different quality-cost trade-offs. We instantiate budget tiers through three complementary tiering strategies: \textit{implementation} tiering (varying the module implementation), \textit{reasoning} tiering (varying inference behavior), and \textit{capacity} tiering (varying the module's model capacity). To navigate these tiered choices, BudgetMem employs a shared lightweight router that performs budget-tier routing as the query is processed: at each module, it selects a tier based on the available context (i.e., the query and intermediate module states). We train the router with reinforcement learning under a cost-aware reward that trades off task performance against memory extraction cost, forming controllable performance–cost behavior.
This design yields a practical and controllable runtime memory system, while also enabling a unified comparison of different budget realization strategies. 


In experiments on LoCoMo, LongMemEval, and HotpotQA, BudgetMem delivers strong gains over competitive baselines in performance-first (high-budget) settings, and exhibits clear performance–cost trade-off curves as budgets tighten. Moreover, our analyses disentangle the relative strengths of different budget tiering strategies, offering insights into which mechanisms provide the best returns under different budget regimes.




Our contributions are summarized as follows:
\begin{itemize}
\item We introduce \textbf{BudgetMem}, a modular runtime agent memory framework that enables explicit performance--cost control for on-demand memory extraction via budget-tiered modules.
\item We propose budget-tier routing, learning a shared lightweight router with reinforcement learning to select budget tiers during extraction and further instantiate three complementary budget realization strategies (i.e., \textit{implementation}, \textit{reasoning}, and \textit{capacity} tiering) within a unified framework.
\item Experiments on LoCoMo, LongMemEval, and HotpotQA demonstrate strong performance and clear performance--cost trade-offs, and our further analyses provide valuable insights into when different strategies deliver the best returns across budget regimes.
\end{itemize}


